\chapter{Critério de parada, liberação do modelo e monitoramento de deriva}
\label{ap:parada-drift}

Este apêndice consolida a política operacional do ciclo de vida do
classificador treinado pelo FALCO, da decisão de parar de rotular à decisão
de retreinar em produção. As duas primeiras seções descrevem mecanismos
implementados na biblioteca \texttt{activelearning}; a terceira é o desenho
recomendado para operação contínua (FlowBuilder), com gatilhos numéricos.

\section{Parada do laço de aprendizado ativo}
\label{ap:parada-laco}

O laço encerra a aquisição de rótulos pelo primeiro de três gatilhos:
\textbf{(i) estagnação na validação}: o Macro F1 em $V$ não melhora mais que
$\varepsilon_{\max}=10^{-3}$ por $p=5$ iterações consecutivas
(Capítulo~\ref{ch:metodo} para o racional das constantes); o conjunto
de teste jamais participa; \textbf{(ii) orçamento}: teto absoluto de rótulos
$B$; \textbf{(iii) exaustão do pool}. O racional do gatilho (i) é
estatístico: o ganho marginal por lote decai ao longo da curva de
aprendizado, enquanto o custo marginal do rótulo é constante; quando o ganho
observável cai abaixo da resolução do próprio conjunto de validação
($\approx 1/\sqrt{n_V}$), rotular mais é pagar para medir ruído. A tolerância
$\varepsilon_{\max}=10^{-3}$ não deriva desse limite amostral: com
$n_V = 2.000$, o limite vale $1/\sqrt{n_V} = 0{,}0224$, e a tolerância é
deliberadamente cerca de vinte e duas vezes menor, de modo que a parada só ocorre
após $p$ iterações sem qualquer ganho registrado, e não por o ganho ser
pequeno demais para se distinguir do ruído. A parada observada no ciclo real,
portanto, não é efeito de uma tolerância estrita: o critério executado é o
mais permissivo dos dois. O gatilho
foi exercitado em duas execuções ponta a ponta da biblioteca, ambas com
\textbf{oráculo LLM real e gratuito} (nemotron via NIM, acurácia
$\approx 78\%$): com orçamento de 1.000, encerrou o ciclo em 991 rótulos
(PVBin) e 982 (SGD); com orçamento de 15.000, encerrou em \textbf{6.009}
rótulos (PVBin) e \textbf{4.742} (SGD), 32--40\% do orçamento, com as
curvas de validação e teste estagnadas em conjunto. Os artefatos estão no
repositório (\texttt{experiments/e5cycle}).

\section{Critério de liberação do modelo}
\label{ap:liberacao}

Recomenda-se liberar o modelo para produção quando três condições valem
simultaneamente: \textbf{(a)} a estagnação foi atingida (o modelo parou de
melhorar com o oráculo disponível); \textbf{(b)} o desempenho de validação
alcança o alvo externo de negócio; na ausência de alvo, o teto empírico do
braço \emph{oráculo-total} (treinar com todo o pool rotulado pelo mesmo
oráculo, ao custo de inferência apenas) indica o máximo que \emph{este}
oráculo permite: estagnar perto dele significa que mais rótulos não ajudam,
apenas um oráculo melhor; \textbf{(c)} o intervalo de confiança de Wilson do
desempenho de validação está inteiro acima do piso aceitável (com
$n_V = 1.000$, a meia-largura é de 2--3 pontos percentuais). Só então o
modelo é avaliado \emph{uma única vez} no conjunto de teste intocado, e esse
é o número reportado.

\section{Monitoramento de deriva e política de retreinamento}
\label{ap:drift}

Textos curtos de varejo derivam continuamente (novos produtos, marcas e
grafias). Três camadas de monitoramento, em ordem de custo:

\begin{enumerate}
  \item \textbf{Deriva de entrada} (sem rótulo, custo zero): fração de tokens
  fora do vocabulário de treino e divergência da distribuição de
  representações (PSI sobre projeções). Gatilho sugerido: taxa de vocabulário
  novo em dobro da taxa do treino, ou PSI $> 0{,}2$.
  \item \textbf{Deriva de confiança} (sem rótulo, custo zero): queda da
  confiança média (ou aumento da entropia média) das predições em produção.
  Gatilho: média móvel de 7 dias abaixo de $\mu - 2\sigma$ do período de
  referência.
  \item \textbf{Verificação amostral com oráculo} (custo quase nulo): acurácia
  numa amostra pequena rotulada pelo LLM Inicial. Com os custos medidos no
  E0, 200 itens semanais custam da ordem de US\$~0,01. Gatilho: acurácia
  amostral abaixo do limite inferior do IC de validação da liberação.
\end{enumerate}

Política: camadas 1--2 disparando antecipam a camada 3; a camada 3
confirmando, retreina-se. O retreino reaproveita todos os rótulos históricos
e roda um ciclo \emph{incremental} de aprendizado ativo apenas sobre o fluxo
recente (o classificador atual seleciona o que não entende; o oráculo LLM
rotula; treina-se do zero com $L_{\text{antigo}} \cup L_{\text{novo}}$;
retreinar do zero, e não continuar do modelo anterior, preserva
comparabilidade e evita a degradação associada a \textit{warm-starting}).
Mesmo sem gatilho, recomenda-se a verificação amostral em cadência mensal:
deriva de catálogo pode ser gradual demais para as camadas 1--2.
